\section{Conditional density and field comparisons}\label{a:sec:conditioning}

We now specialize to one nucleus at the origin and use
$V(y)=Z/|y|$, $d_y=|y|$, and the atomic distance scale
$a_y=d_y/L$ with $L=10^5$. The comparisons in this section allow every
real $Z\ge1$; the neutral-sector conclusions below use integer charges.

Fix $0\le D_0<\infty$ and a normalized ground state $\Psi$ in a sector
with $n\le3Z$ and $E_n\le E+D_0$. The ground-state assumption in
this section is used in the multiplier identity
\eqref{eq:sector-multiplier}; the patch comparisons themselves apply
to the resulting states with an energy offset. All probabilities
below initially refer to the raw position-and-spin law of $\Psi$,
enlarged by the explicitly specified independent random variables.
The common likelihood and event-cost estimates of
Section~\ref{sec:observations} now give a simultaneous comparison throughout
an atomic annulus. This strengthens the pointwise molecular comparison by
using the absence of nuclear singularities from these annuli.

\subsection{Nested observations and master smoothing}

Choose $0<\eps<1$ and $0<s<1$, and suppose $Z$ is large enough that
$2r_0\le s$. Set
\begin{equation}\label{a:eq:cond-observations}
 r_0=\eps Z^{-1/3},\qquad r_j=2^jr_0,\qquad
 m=\max\{j:r_j\le s\},\qquad \ell_j=r_j^{1.01}\quad(j<m).
\end{equation}
Thus $s/2<r_m\le s$. At scale $j<m$ observe the positions
$x_i+\ell_j U_{ji}$ and then forget the particle labels separately
in each array. All Cartesian coordinates of all $U_{ji}$, at all
scales, are independent with density
\[
 \zeta(u)=c_\zeta\exp\!\left(-\frac1{1-u^2}\right)
                  \1_{\{|u|<1\}}.
\]
Represent the unordered array by its lexicographically sorted tuple
$Y_j$. Ties have probability zero. Write
$\mathcal F_j=\sigma(Y_j,\ldots,Y_{m-1})$ and let
$\mathcal F_m$ be trivial. These sigma-fields decrease with $j$.

Fix $w=10^{-5}$ and $0<c_1<(10L)^{-1}$. With the same smooth radial
packet $g$ as in Section~\ref{sec:localization}, define
\begin{equation}\label{a:eq:cond-master}
 \begin{gathered}
 D_x^0=\max(d_x,r_0),\qquad
 t_x=c_1D_x^0\min(D_x^0,s)^w,\qquad
 \mathcal K_x(y)=g_{t_x}(y-x)^2,\\
 \mu_j(y)=\E\!\left[\sum_{i=1}^n\mathcal K_{x_i}(y)
                    \,\middle|\,\mathcal F_j\right],\qquad
 H_j=V-K*\mu_j,\qquad
 \mathcal R f(y)=\int\mathcal K_z(y)f(z)\,\dd z.
 \end{gathered}
\end{equation}
The lower clamp in $D_x^0$ is retained also for particles close to
the nucleus; the upper clamp is only on the factor raised to $w$.
In particular, $t_x\ge t_{\min}:=c_1r_0^{1+w}>0$ and $t$ is
Lipschitz with constant at most $c_1(1+w)s^w$.


Use the likelihood-ratio versions of Lemma~\ref{lem:obs-data} for these
arrays and kernels. They give joint measurability, fixed-system derivative
bounds, $\mu_j\ge0$, $\int\mu_j=n$, and
\begin{equation}\label{a:eq:cond-tower}
 \E[\mu_j\mid\mathcal F_{j+1}]=\mu_{j+1},\qquad
 \E[H_j\mid\mathcal F_{j+1}]=H_{j+1}.
\end{equation}
The offset $H_j-V$ is locally Lipschitz and
\begin{equation}\label{a:eq:cond-poisson}
 \Delta H_j=-4\pi\nu+4\pi\mu_j.
\end{equation}
Lemma~\ref{lem:obs-tilt} applies with baseline $E_n$ and offset $D_0$.
In particular, an event of probability $p>0$ at index $j$ has tilted
energy at most $E+D_0+C\ell_j^{-2}\log^5(e/p)$. We retain this full
baseline dependence throughout the argument.

\subsection{The simultaneous inverse comparison}

\begin{proposition}[Conditional density determines the field]
\label{a:prop:cond-comparison}
There is a universal constant $C_{\rm cap}$ with the following
property. Fix $L_1\ge2$, $0<h_l<h_h<\infty$, and a sufficiently
small $\xi>0$; it suffices to require $\xi<h_l/16$.
There exists
$s_*=s_*(D_0,L_1,h_l,h_h,\xi)>0$ such that, whenever
$0<s\le s_*$, at each $j<m$, with $r=r_j$, there is an event
$\mathcal G_j\in\mathcal F_j$ satisfying
$\Prob(\mathcal G_j^c)\le Cr^{25}$ on which, simultaneously for
$r\le d_y\le4L_1r$ and $h_l\le h\le h_h$,
\begin{equation}\label{a:eq:cond-inverse}
 \begin{aligned}
 d_y^6\mu_j(y)>kh^{3/2}
       &\quad\Longrightarrow\quad d_y^4H_j(y)\ge h-\xi,\\
 d_y^6\mu_j(y)<kh^{3/2}
       &\quad\Longrightarrow\quad d_y^4H_j(y)\le h+\xi.
 \end{aligned}
\end{equation}
On the same event, $H_j(y)\le C_{\rm cap}d_y^{-4}$ throughout
that band. The constant $C$ may depend on the fixed parameters,
whereas $C_{\rm cap}$ does not depend on $D_0,L_1,h_l,h_h,\xi$.
All estimates are uniform in $Z$, the ground state, $j$, and $m$.
\end{proposition}

\begin{proof}
We first control the observed count and the positive field throughout
the band. At a fixed point and height, a hypothetical failure then
defines a tilted state. Its fresh patch compares the observed density
with the patch field, including the contribution of rare negative
fields; the two expectation towers turn this into a contradiction.
Finally, spatial and height nets give the simultaneous statement.
Every occurrence of $o(1)$ below tends to zero as $s\downarrow0$,
uniformly for $r_0\le r=r_j\le s$ and the stated bands, after the
fixed parameters in the proposition have been chosen.

\emph{Kernel geometry.}
Put $a=a_y=d_y/L$, $t=t_y$ and $\tau=t/a$.
If $\mathcal K_z(y)\ne0$, the Lipschitz bound on $t$ gives
$|z-y|\le t_z\le t_y+Cs^w|z-y|$, hence $|z-y|\le2t_y$.
On this support, $t_z/t_y=1+O(s^w)$. Comparison with the fixed-width
kernel $g_t(y-z)^2$ therefore gives
\begin{equation}\label{a:eq:cond-kernel}
 \int\mathcal K_z(y)\,\dd z=1+O(s^w),\qquad
 |\mathcal K_z(y)|\le Ct^{-3},\qquad
 |\nabla_y\mathcal K_z(y)|+|\nabla_z\mathcal K_z(y)|\le Ct^{-4}.
\end{equation}
The center derivative is a weak derivative where $t_z$ has a
corner; its stated bound also gives the needed Lipschitz estimate.
For $r/2\le d_y\le8L_1r$ one has $a\asymp_{L_1}r$ and
\begin{equation}\label{a:eq:cond-tau}
 c_{L_1}a^w\le\tau\le Cs^w.
\end{equation}
For example, when $d_y\ge r_0$, the formula is
$\tau=c_1L\min(d_y,s)^w$; on the indicated band
$\min(d_y,s)\ge c_{L_1}d_y$. The case $d_y<r_0$ is comparable,
since $d_y\ge r_0/2$.

\emph{An observed-count event.}
Let $M_j$ be the number of entries of $Y_j$ with radii in
$[r/4,12L_1r]$. For a sufficiently large fixed $C_{L_1}$,
\begin{equation}\label{a:eq:cond-count-event}
 \Prob\{M_j>C_{L_1}r^{-3}\}\le r^{40}.
\end{equation}
Indeed, if the event had probability $p>r^{40}$, its tilted raw
state from Lemma~\ref{lem:obs-tilt} would have offset
\[
 \delta\le D_0+Cr^{-2.02}\log^5(e/r^{40}).
\]
Each configuration compatible with the event has more than
$C_{L_1}r^{-3}$ raw points in $[r/8,13L_1r]$, because observation
displacements are at most $\sqrt3\ell_j=o(r)$. But
\eqref{a:eq:screen-annular} bounds its expected count by
$C'_{L_1}(r^{-3}+\sqrt{\delta r})\le2C'_{L_1}r^{-3}$ for small
$s$, a contradiction if $C_{L_1}>2C'_{L_1}$. Call the complementary
event in \eqref{a:eq:cond-count-event} $\mathcal C_j$.
For data in $\mathcal C_j$, every compatible configuration has at
most $C_{L_1}r^{-3}$ centers relevant to the master kernels on
$r/2\le d_y\le8L_1r$. Conditional averaging and
\eqref{a:eq:cond-kernel} imply there
\begin{equation}\label{a:eq:cond-count-regularity}
 \mu_j(y)\le Cr^{-6-3w},\qquad
 |\nabla\mu_j(y)|\le Cr^{-7-4w}.
\end{equation}

\emph{A universal spatial field cap.}
Fix $y$ in the same enlarged band and tilt an event
$A'=\{H_j(y)>b\,d_y^{-4}\}$ of positive probability $p$.
The original-law tower gives
\[
 \E[H_j(y)\mid A']
  =V(y)-\E_{\Psi_{A'}}\sum_i(K*\mathcal K_{x_i})(y)
  \le\E_{\Psi_{A'}}J_y.
\]
To justify the last inequality, a master smear whose potential at
$y$ differs from the point potential has
$|x_i-y|<t_{x_i}$ and hence $|x_i-y|\le2t_y<Aa_y$ for small
$s$. All terms outside that exclusion therefore have their exact
point potential, and dropping the remaining electron terms gives
$J_y$. By \eqref{a:eq:screen-local} without an initial cut and
Lemma~\ref{lem:obs-tilt},
\begin{equation}\label{a:eq:cond-positive-tail}
 b\le C_2+C_{L_1}r^{2.49}\log^{5/2}(e/p).
\end{equation}
Here $C_2$ is universal: the local field term in these units is
$C d_y^4a_y^{-4}=CL^4$, and the contribution
$C d_y^{7/2}\sqrt{D_0}$ is at most one after decreasing $s$.
Solving \eqref{a:eq:cond-positive-tail} for $p$ shows that for
$b\ge2C_2$,
\[
 \Prob\{d_y^4H_j(y)>b\}
   \le e\exp\{-c_{L_1}(b-C_2)^{2/5}r^{-0.996}\}.
\]
Tail integration consequently gives, for each fixed $M>0$,
\[
 \E(d_y^4H_j(y)-2C_2)_+\le C_{M,L_1}r^M.
\]
Subharmonicity permits a simultaneous conclusion without a
$Z$-dependent derivative estimate. Apply the ball submean inequality
with radius $d_y/10$ for $3r/4\le d_y\le6L_1r$. The averaging ball
is contained in $\{r/2\le d_z\le8L_1r\}$ and $d_z\asymp d_y$
there, so
\[
 d_y^4H_j(y)\le C_3+Cr^{-3}
   \int_{\{r/2\le d_z\le8L_1r\}}
                (d_z^4H_j(z)-2C_2)_+\,\dd z.
\]
The constant $C_3$ is universal. The integral term has expectation
$O_{M,L_1}(r^M)$, since its integration volume is $O_{L_1}(r^3)$.
Markov's inequality proves that, outside an event of probability
$C_{M,L_1}r^M$,
\begin{equation}\label{a:eq:cond-cap}
 H_j(y)\le C_{\rm cap}d_y^{-4}
       \quad(3r/4\le d_y\le6L_1r),\qquad C_{\rm cap}=C_3+1.
\end{equation}
This establishes the claimed order of dependence of the cap.

\emph{A fixed spatial point and height.}
Fix now $r\le d_y\le4L_1r$ and
$h\in[h_l/2,2h_h]$. Consider either failure of
\eqref{a:eq:cond-inverse} with tolerance $\xi/4$, intersected with
$\mathcal C_j$. If its probability were at least $r^{42}$, call
that event $A'$ and write $Q=\Prob(\,\cdot\mid A')$.
The raw marginal is the state $\Psi_{A'}$ of
Lemma~\ref{lem:obs-tilt}, with
\begin{equation}\label{eq:obs-tilted-offset}
 \delta\le D_0+Cr^{-2.02}\log^5(e/r^{42})
                    \le a^{-7+.02}
\end{equation}
for small $s$. Apply the refined patch
\eqref{a:eq:screen-refined} to this raw state at $y$, with
$\beta=a^{.2}$ and $b=\beta a$. Extend $Q$ by fresh cut labels
sampled conditional on the raw positions, exactly as in that patch.
Let $X$ be these patch data, let $\rho_X^c$ be its conditional core
density, and let $\rho,W,\sigma$ be its minimizing density,
screened field, and fine-smear density. Thus
\[
 W=V-K*\rho_X^c-K*\rho,\qquad
 \rho=kW_+^{3/2}\ \hbox{on the patch},\qquad
 \E_Q\left[D(\sigma-\rho)+\int W_-\sigma\right]
       \le Ca^{-7+.02}.
\]
The roles of these densities and their laws are distinct. The pair
$\mu_j,H_j$ keeps its original-observation definition in
\eqref{a:eq:cond-master}. Under the event law $Q$, the fresh patch gives
the conditional core density $\rho_X^c$ and its field
$\Phi_X=V-K*\rho_X^c$; $\sigma$ is the fine smear of retained out
particles, whereas $\rho$ is the auxiliary minimizing density and
$W=\Phi_X-K*\rho$. The field comparison is made after taking
$Q$-expectations in \eqref{a:eq:cond-transfer} below.

Each compatible raw configuration matches the observed array
$Y_j$ after a permutation, with displacement at most
$\sqrt3\ell_j$ per position. Both configurations and their
conditional averages at those same data therefore differ, when
tested against $x\mapsto\mathcal K_x(y)$, by at most
$Cr^{-3}\ell_jt^{-4}$. This follows from the count bound and the
center Lipschitz estimate in \eqref{a:eq:cond-kernel}; only entries
within $2t+O(\ell_j)$ of $y$ can contribute. All such raw particles
are retained by the fresh patch. Its radial fine smear moves a
retained point by at most $\sqrt3b$, adding at most
$Cr^{-3}bt^{-4}$ to the error. Consequently, on every compatible
path, including paths where the Coulomb error is large,
\begin{equation}\label{a:eq:cond-matching}
 |\mathcal R\sigma(y)-\mu_j(y)|
   \le Cr^{-3}(\ell_j+b)t^{-4}
   \le Ca^{-6}\big(a^{.01}\tau^{-4}+a^{.2}\tau^{-4}\big)
   =o(a^{-6}).
\end{equation}
In the comparison with the conditional average, one may first
compare each raw sum with $\sum_i\mathcal K_{Y_{j,i}}(y)$, then
average this same bound using the original posterior law. The
array is fixed in both comparisons, so no equality of posterior
laws is being used.

Let $G=\{D(\sigma-\rho)\le a^{-7+.01}\}$. Then
$Q(G^c)\le Ca^{.01}$. The test function
$z\mapsto\mathcal K_z(y)$ has support in $B(y,2t)$ and gradient
norm at most $Ct^{-5/2}$. Coulomb duality gives on $G$
\begin{equation}\label{a:eq:cond-smoothed-duality}
 |\mathcal R(\sigma-\rho)(y)|
  \le Ca^{-7/2+.005}t^{-5/2}
  =Ca^{-6}a^{.005}\tau^{-5/2}=o(a^{-6}).
\end{equation}
The positivity of the two powers used so far is explicit:
$.01-4w=.00996$ and $.005-(5/2)w=.004975$.

On $G$, a high-density antecedent forces
\begin{equation}\label{a:eq:cond-patch-high}
 W(y)\ge(h-\xi/16)d_y^{-4}.
\end{equation}
For if the reverse strict inequality held, the oscillation bound
\eqref{a:eq:screen-oscillation} on $B(y,2t)$ would give
$W(z)\le(h-\xi/32)d_y^{-4}$ there for sufficiently small $s$.
This remains uniform when $W(y)$ is arbitrarily negative:
the upper estimate is $W(y)+C\tau(a^{-4}+W(y)_-)$, and
$v\mapsto v+C\tau v_-$ is increasing for $C\tau<1$.
Using the patch equation and
$\int\mathcal K_z(y)\,\dd z=1+O(s^w)$ would then give
\[
 \mathcal R\rho(y)
   \le k(h-\xi/32)^{3/2}d_y^{-6}(1+O(s^w))
   <kh^{3/2}d_y^{-6}-c_{h_l,h_h,\xi}a^{-6},
\]
contradicting \eqref{a:eq:cond-matching}--\eqref{a:eq:cond-smoothed-duality}.
Similarly, on $G$ a low-density antecedent forces
\begin{equation}\label{a:eq:cond-patch-low}
 W(y)\le(h+\xi/16)d_y^{-4}.
\end{equation}
Indeed the contrary is positive, so the lower oscillation estimate
gives $W(z)\ge(h+\xi/32)d_y^{-4}$ on the contributing ball; its
reaction density contradicts the low antecedent with the same
strict margin. Constants from $d_y=La$ are fixed before $s$ is
decreased.

\emph{The negative field on the exceptional patch event.}
The upper bound \eqref{a:eq:cond-patch-low} passes to expectation
with an $o(a^{-4})$ error, by the patch bound $W(y)\le Ca^{-4}$
and $Q(G^c)\le Ca^{.01}$. To pass the lower bound to expectation
we need more than this small probability. Under the high-density
antecedent, \eqref{a:eq:cond-matching} holds on all paths and gives
$\mathcal R\sigma(y)\ge ca^{-6}$, since $h\ge h_l/2>0$.
The support and upper bound of the master kernel imply the
deterministic mass bound
\[
 M_\sigma:=\int_{B(y,2t)}\sigma\ge ca^{-6}t^3.
\]
For $Y=W(y)_-$, \eqref{a:eq:screen-oscillation} gives throughout this
ball
$W(z)_-\ge(1-C\tau)Y-C\tau a^{-4}$.
Taking $C\tau\le1/2$, integrating against $\sigma$, and using the
mass bound yields
\[
 Y\le C(a^{-6}t^3)^{-1}\int W_-\sigma+C\tau a^{-4}.
\]
Multiply by $\1_{G^c}$ and use the refined patch budget under the
same law $Q$. We obtain
\begin{equation}\label{a:eq:cond-negative-exception}
 \E_Q[\1_{G^c}W(y)_-]
  \le Ca^{-4}\big(a^{.02}\tau^{-3}+\tau\big)=o(a^{-4}),
\end{equation}
because $.02-3w=.01997>0$. Thus the high and low antecedents give,
respectively,
\begin{equation}\label{a:eq:cond-patch-means}
 \E_QW(y)\ge(h-\xi/16)d_y^{-4}-o(a^{-4}),\qquad
 \E_QW(y)\le(h+\xi/16)d_y^{-4}+o(a^{-4}).
\end{equation}


\emph{Transfer between the two conditioning procedures.}
Apply Lemma~\ref{lem:obs-transfer} with $V_s=V$, $V_c=0$ and
$h_X=W$. Its two towers are taken under the original observation law
and the fresh tilted patch law, respectively. The refined patch provides
all its inputs: the raw density bound and Coulomb budget are
\eqref{a:eq:screen-refined}, the deleted count bound is
\eqref{m:eq:screen-deleted-shell}, and $\rho\le Ca^{-6}$ on the
contributing ball by \eqref{a:eq:screen-interior}. Therefore
\begin{equation}\label{a:eq:cond-transfer}
 |\E_Q(H_j-W)(y)|
 \le Ca^{-4}\{\tau^{1/5}+\beta^{1/5}+\sqrt\beta
             +a^{.01}\tau^{-1/2}+\tau^{1/2}\}=o(a^{-4}).
\end{equation}
Here $\beta=a^{.2}$ and \eqref{a:eq:cond-tau} makes every positive-power
error vanish. The source-free potential remainder actually has the sharper
power $\tau^2$, so the common upper bound is sufficient. No equality of
posterior laws is used in this transfer.

On a high-failure event, its defining inequality is
$d_y^4H_j(y)<h-\xi/4$, whereas
\eqref{a:eq:cond-patch-means} and \eqref{a:eq:cond-transfer} give
$d_y^4\E_QH_j(y)\ge h-\xi/16-o(1)$, a contradiction for small
$s$. On a low-failure event the inequalities are reversed and
give the same contradiction. Thus each such failure intersected
with $\mathcal C_j$ has probability less than $r^{42}$. Including
\eqref{a:eq:cond-count-event}, the two implications at a fixed
point and height, with tolerance $\xi/4$, fail with probability
at most $Cr^{40}$.

\emph{Spatial and height nets.}
Choose a net of the closed band $\{r\le d_y\le4L_1r\}$ with
spacing at most $r^2$ and at most $C_{L_1}r^{-3}$ points. For
example, select one band point in each intersecting cube of side
$r^2/2$; the cube count follows from the volume of its $O(r^2)$
enlargement. Choose a finite height net in $[h_l/2,2h_h]$ with
mesh at most $\xi/64$, including its endpoints. Exclude every
point-and-height failure, the count failure, and the cap failure
with $M=40$ in \eqref{a:eq:cond-cap}. A union bound gives exceptional
probability $Cr^{37}$, and hence at most $Cr^{25}$ for $r<1$.

On the retained event, \eqref{a:eq:cond-count-regularity} shows that
$d_y^6\mu_j(y)$ varies by at most $Cr^{1-4w}=o(1)$ between a
band point and a net point within $Cr^2$. To control the field,
work in a ball of radius $cr$ around a band point, contained in
the enlarged cap band. On this ball $B$, let
$Q=K*(\1_B\mu_j)$. Then $H_j+Q$ is harmonic on $B$.
The value and gradient of $Q$ are bounded by
$Cr^{-4-3w}$ and $Cr^{-5-3w}$, respectively, so $H_j+Q$ is
bounded above by $Cr^{-4-3w}$. Its nonnegative distance from that
upper bound has, by the Poisson formula on a smaller concentric
ball, oscillation bounded by its central value times the relative
distance. It follows that
\begin{equation}\label{a:eq:cond-field-oscillation}
 |H_j(z)-H_j(y)|
   \le Cr\big(r^{-4-3w}+H_j(y)_-\big)
       \quad(|z-y|\le Cr^2).
\end{equation}
This estimate also controls the change of the normalized field
$d_y^4H_j(y)$, since $d_z/d_y=1+O(r)$. In particular, if its
value at $y$ is below $h-\xi$, its value at $z$ is below
$h-3\xi/4$ for small $s$. Arbitrarily negative values cause no
problem: the relevant upper estimate has the form
$v+Crv_-+O(r^{1-3w})$, increasing in $v$ when $Cr<1$.
If the value at $y$ exceeds $h+\xi$, it is positive and the cap
and \eqref{a:eq:cond-field-oscillation} make its value at $z$
greater than $h+3\xi/4$.

Suppose now a high-density implication failed at some $(y,h)$.
At a nearby spatial net point choose a height-net value
$h'\in[h-\xi/4,h-\xi/8]$. The positive lower height bound and
the $o(1)$ density variation preserve the high antecedent at
$h'$, while the field is below $h-3\xi/4<h'-\xi/4$.
This is an excluded net failure. For a low-density failure choose
$h'\in[h+\xi/8,h+\xi/4]$; its antecedent is preserved and the
field exceeds $h+3\xi/4>h'+\xi/4$. The height intervals lie
inside $[h_l/2,2h_h]$ by the stipulated smallness of $\xi$.
This proves simultaneity and finishes the proposition.
\end{proof}

We record also an unconditional estimate used to pay for a failed
band. For fixed $L_1$ and $r\le d_y\le4L_1r$, the support and
size of the master kernel, conditional Jensen, and
\eqref{a:eq:screen-annular} in the original state give
\begin{equation}\label{a:eq:cond-density-l2}
 \|\mu_j(y)\|_2
  \le Ct_y^{-3}
    \left\|\#\{i:r/2<d_{x_i}<8L_1r\}\right\|_2
  \le Ct_y^{-3}r^{-3}\le Cr^{-6-3w}.
\end{equation}
The middle inequality uses the offset $D_0$ and sufficiently small
$s$; its constant is uniform in the state and scale.

The order of choices is relevant. The fixed screening constants,
$g,c_1,w$, and the universal cap $C_{\rm cap}$ come first.
One may then choose $L_1,h_l,h_h,\xi$ and a bounded offset $D_0$,
and finally make $s$ small enough for every estimate above.
Thereafter any sufficiently large $Z$ with $2r_0\le s$ is
permitted. In particular no step requires $D_0=0$, and the
number of dyadic scales introduces no loss in the constants.
